\chapter{Complément : fonder les outils d'analyse}
\section*{Exercice 25.1 -- Une limite quantifiée}
Donner un $\delta$ explicite prouvant que $x^2\to4$ lorsque $x\to2$.

\section*{Exercice 25.2 -- Hypothèses des bornes atteintes}
Donner un exemple d'une fonction continue et bornée sur $]0,1[$ qui n'atteint ni sa borne supérieure ni sa borne inférieure.

\section*{Exercice 25.3 -- Valeurs intermédiaires}
Prouver que $x^3+x-1$ possède exactement une racine dans $[0,1]$.

\section*{Exercice 25.4 -- Dichotomie}
Une racine est cherchée par dichotomie dans $[0,1]$. Combien de divisions suffisent pour que le milieu de l'intervalle final ait une erreur au plus $10^{-3}$ ?

\section*{Exercice 25.5 -- Intégrale par sommes}
Retrouver $\int_0^1t\,dt$ par des sommes à gauche et à droite sur $n$ intervalles égaux.

\section*{Exercice 25.6 -- Cauchy}
Si $|u_{n+1}-u_n|\le3^{-n}$ pour tout $n\ge0$, prouver que $(u_n)$ converge dans $\R$.
